\documentclass[11pt]{article}

\usepackage[margin=1in]{geometry}
\usepackage{amsmath, amssymb, amsthm}
\usepackage{bm}
\usepackage{booktabs}
\usepackage{enumitem}
\usepackage{hyperref}

\hypersetup{colorlinks=true, linkcolor=blue, citecolor=blue, urlcolor=blue}

% --- Notation shortcuts ---
\newcommand{\R}{\mathbb{R}}
\newcommand{\E}{\mathbb{E}}
\newcommand{\norm}[1]{\left\lVert #1 \right\rVert}
\newcommand{\fro}[1]{\norm{#1}_{F}}
\newcommand{\twonorm}[1]{\norm{#1}_{2}}
\newcommand{\T}{^{\!\top}}
\newcommand{\range}{\operatorname{span}}
\newcommand{\rank}{\operatorname{rank}}
\newcommand{\diag}{\operatorname{diag}}
\newcommand{\Cov}{\operatorname{Cov}}
\newcommand{\orth}{\operatorname{orth}}
\DeclareMathOperator*{\argmin}{arg\,min}

\title{\textbf{Evaluation Metrics for the Descriptor Diffusion Factor Model}\\[4pt]
\large Mathematical Definitions}
\author{Diffusion Learn Factors --- branch \texttt{subspace-recovery-metric}}
\date{\today}

\begin{document}
\maketitle

\begin{abstract}
This note states, precisely and in self-contained notation, every quantitative
metric computed in the descriptor diffusion factor model. We group them into
(i)~\emph{subspace recovery} metrics that compare the learned descriptor-to-factor
map $A$ against the ground-truth map $T$ in a rotation- and scaling-invariant way,
(ii)~\emph{distribution-matching} sanity metrics comparing real and generated
returns through their first two moments, and (iii)~\emph{ground-truth subspace and
eigenvalue} errors that benchmark both the training sample and the generated sample
against the analytically known return covariance. Each definition is paired with a
one-line interpretation.
\end{abstract}

%==================================================================
\section{Setup and notation}
%==================================================================

\paragraph{Dimensions.}
$d$ = number of assets, $m$ = number of descriptors, $k$ = number of latent
factors, $n$ = number of samples (periods). Throughout, $\twonorm{\cdot}$ is the
Euclidean norm, $\fro{\cdot}$ the Frobenius norm, and for $M\in\R^{a\times b}$
the singular values are written $\sigma_1(M)\ge\sigma_2(M)\ge\cdots\ge 0$.

\paragraph{Generative model (residual-free).}
For period $i$ we are given a descriptor matrix $U_i\in\R^{d\times m}$ and a
\emph{time-invariant} ground-truth factor map $T\in\R^{m\times k}$. Latent factors
are standard Gaussian, $F_i\sim\mathcal N(0,I_k)$, and returns are generated noiselessly as
\begin{equation}
  r_i \;=\; U_i\,T\,F_i \;\in\; \R^{d}.
  \label{eq:genmodel}
\end{equation}
In the \emph{fixed} regime $U_i\equiv U$ is constant across periods, so the
asset loadings $\beta = U T\in\R^{d\times k}$ are time invariant and
$r_i=\beta F_i$. The identifiable object is the \emph{loading subspace}
$\range(\beta)$; both $T$ and $A$ (below) are only defined up to a $k\times k$
change of basis.

\paragraph{Score decomposition (the model).}
Let $Q_i\in\R^{d\times m}$ be the orthonormal factor of the thin QR decomposition
of $U_i$ (so $Q_i\T Q_i = I_m$ and $\range(Q_i)=\range(U_i)$), and write the
projected state $z = Q_i\T r$. With a learned map $A\in\R^{m\times k}$ and a
network $g_\zeta(\cdot)\in\R^k$, the score estimator is
\begin{equation}
  s_\theta(r,t,i) \;=\; -\frac{1}{h_t}\,r \;+\; \frac{\alpha_t}{h_t}\,U_i\,A\;
  g_\zeta\!\bigl(Q_i\T r,\,t,\,i\bigr),
  \label{eq:score}
\end{equation}
where $\alpha_t=\sqrt{\bar\alpha_t}$ and $h_t = 1-\bar\alpha_t$ are the standard
DDPM forward-process coefficients ($x_t=\alpha_t x_0+\sqrt{h_t}\,\varepsilon$,
$\varepsilon\sim\mathcal N(0,I_d)$, $\bar\alpha_t=\prod_{s\le t}\alpha_s$).
The network is trained on the noise-prediction reparameterisation
$\varepsilon_\theta = -\sqrt{h_t}\,s_\theta$ with the DDPM objective
$\mathcal L = \E_{t,x_0,\varepsilon}\fro{\varepsilon_\theta(x_t,t)-\varepsilon}^2$.
The matrix $A$ is the model's estimate of $T$; comparing $A$ to $T$ is the central
identifiability question, addressed in Section~\ref{sec:subspace}.

%==================================================================
\section{Subspace recovery metrics: \texorpdfstring{$A$}{A} versus \texorpdfstring{$T$}{T}}
\label{sec:subspace}
%==================================================================

Because the factors are only identifiable up to a $k\times k$ invertible map,
$A$ and $T$ should \emph{not} be compared entrywise. We instead compare their
column spaces, $\range(A)$ and $\range(T)$, with rotation- and scaling-invariant
quantities.

\subsection{Orthonormal bases}
For a matrix $M\in\R^{m\times k}$ with thin SVD $M = U_M\,\Sigma_M\,V_M\T$, define
its numerical rank and an orthonormal basis of its column space by
\begin{equation}
  r(M) = \#\bigl\{\, j : \sigma_j(M) > \tau\,\max(\sigma_1(M),1) \,\bigr\},
  \qquad
  Q_M = \orth(M) = \bigl[\,u_1\ \cdots\ u_{r(M)}\,\bigr],
\end{equation}
i.e.\ the first $r(M)$ left singular vectors, with tolerance $\tau = 10^{-9}$.
We write
\begin{equation}
  Q_A = \orth(A)\in\R^{m\times r_A},\quad
  Q_T = \orth(T)\in\R^{m\times r_T},\quad
  k_\star = \min(r_A, r_T),
\end{equation}
so that $Q_A\T Q_A = I_{r_A}$ and $Q_T\T Q_T = I_{r_T}$.

\subsection{Principal angles}
The principal angles $0\le\theta_1\le\cdots\le\theta_{k_\star}\le\tfrac{\pi}{2}$
between $\range(A)$ and $\range(T)$ have cosines equal to the singular values of
$Q_A\T Q_T$:
\begin{equation}
  \cos\theta_j \;=\; \operatorname{clip}_{[-1,1]}\!\bigl(\sigma_j(Q_A\T Q_T)\bigr),
  \qquad
  \theta_j = \arccos(\cos\theta_j), \quad j = 1,\dots,k_\star.
  \label{eq:angles}
\end{equation}
We report (in degrees) the full vector $(\theta_1,\dots,\theta_{k_\star})$ together
with its mean and maximum:
\begin{equation}
  \overline{\theta} = \frac{1}{k_\star}\sum_{j=1}^{k_\star}\theta_j,
  \qquad
  \theta_{\max} = \max_{j}\theta_j .
\end{equation}
\emph{Interpretation:} $\theta_j=0$ for all $j$ means the subspaces coincide;
$\theta_{\max}$ near $90^\circ$ signals a direction of $\range(T)$ not captured by
$\range(A)$.

\subsection{Grassmann (geodesic) distance}
The geodesic distance on the Grassmann manifold is the $\ell_2$ norm of the angle
vector:
\begin{equation}
  d_{\mathrm{G}}(A,T) \;=\; \twonorm{\bm\theta} \;=\;
  \Bigl(\textstyle\sum_{j=1}^{k_\star}\theta_j^2\Bigr)^{1/2}.
  \label{eq:grassmann}
\end{equation}
\emph{Interpretation:} a single scalar summary of subspace disagreement; $0$ iff
the subspaces coincide.

\subsection{Projector (Frobenius) distance}
Let $P_A = Q_A Q_A\T$ and $P_T = Q_T Q_T\T$ be the orthogonal projectors onto the
two subspaces. The projector distance and its normalisation are
\begin{equation}
  \Delta_P \;=\; \fro{P_A - P_T},
  \qquad
  \widehat{\Delta}_P \;=\; \frac{\Delta_P}{\sqrt{2\,k_\star}} .
  \label{eq:projector}
\end{equation}
For equal-dimensional subspaces ($r_A=r_T=k_\star$) this is exactly tied to the
principal angles via
\begin{equation}
  \fro{P_A - P_T} \;=\; \sqrt{2}\,\twonorm{\sin\bm\theta}
  \;=\; \sqrt{2}\Bigl(\textstyle\sum_{j}\sin^2\theta_j\Bigr)^{1/2},
\end{equation}
and the maximum possible value over two $k_\star$-dimensional subspaces is
$\sqrt{2k_\star}$, so $\widehat\Delta_P\in[0,1]$.
\emph{Interpretation:} a basis-free distance between the subspaces; $0$ = perfect
overlap, $1$ = orthogonal subspaces.

\subsection{Orthogonal Procrustes residual}
We align the bases by the optimal orthogonal map
$R^\star = \argmin_{R\T R = I}\fro{Q_A R - Q_T}$. With the SVD
$Q_A\T Q_T = W\,\Sigma\,Z\T$ the solution is $R^\star = W Z\T$, and we report the
relative residual
\begin{equation}
  \rho \;=\; \frac{\fro{Q_A R^\star - Q_T}}{\max\bigl(\fro{Q_T},\,10^{-12}\bigr)},
  \qquad \fro{Q_T}=\sqrt{r_T}.
  \label{eq:procrustes}
\end{equation}
\emph{Interpretation:} the fraction of $Q_T$'s energy that the best rotation of
$Q_A$ fails to reproduce; $\rho=0$ means $\range(A)$ contains $\range(T)$ exactly.

\subsection{Singular-value spectra (reference only)}
For completeness we also record the raw spectra
$\sigma(A)=(\sigma_1(A),\dots)$ and $\sigma(T)=(\sigma_1(T),\dots)$.
These are \emph{not} rotation invariant and are blind to the column space; they are
reported only to diagnose scaling/conditioning, never as a recovery criterion.

%==================================================================
\section{Distribution-matching (moment) metrics}
\label{sec:moments}
%==================================================================

Let $\{r^{(s)}\}_{s=1}^{n}\subset\R^{d}$ be real (training) returns and
$\{g^{(s)}\}_{s=1}^{n}\subset\R^{d}$ generated returns. Define the empirical means
and (unbiased, $n-1$ normalised) covariances
\begin{equation}
  \mu_{\mathrm{r}} = \frac1n\sum_{s} r^{(s)},\quad
  \mu_{\mathrm{g}} = \frac1n\sum_{s} g^{(s)},\quad
  C_{\mathrm{r}} = \frac{1}{n-1}\sum_{s}(r^{(s)}-\mu_{\mathrm r})(r^{(s)}-\mu_{\mathrm r})\T,
\end{equation}
and $C_{\mathrm g}$ analogously. The reported metrics are:
\begin{align}
  \text{mean norms:}\quad
    &\twonorm{\mu_{\mathrm r}},\quad \twonorm{\mu_{\mathrm g}},\quad
     \twonorm{\mu_{\mathrm r}-\mu_{\mathrm g}}, \\[2pt]
  \text{covariance norms:}\quad
    &\fro{C_{\mathrm r}},\quad \fro{C_{\mathrm g}},\quad
     \fro{C_{\mathrm r}-C_{\mathrm g}}, \\[2pt]
  \text{covariance ratio:}\quad
    &\kappa_{\mathrm{cov}} = \frac{\fro{C_{\mathrm g}}}{\fro{C_{\mathrm r}}}, \\[2pt]
  \text{relative covariance error:}\quad
    &\epsilon_{\mathrm{cov}} = \frac{\fro{C_{\mathrm r}-C_{\mathrm g}}}{\fro{C_{\mathrm r}}}, \\[2pt]
  \text{std ratio:}\quad
    &\kappa_{\mathrm{std}} = \frac{\frac1d\sum_{j=1}^{d}\widehat\sigma_j(g)}
                                  {\frac1d\sum_{j=1}^{d}\widehat\sigma_j(r)},
\end{align}
where $\widehat\sigma_j(\cdot)$ is the per-asset sample standard deviation
(population normalisation, $n$). \emph{Interpretation:} $\twonorm{\mu_{\mathrm r}-\mu_{\mathrm g}}$
and $\epsilon_{\mathrm{cov}}$ should be small; $\kappa_{\mathrm{cov}},\kappa_{\mathrm{std}}\to 1$
indicate the generator matches the real second moments rather than under- or
over-dispersing.

%==================================================================
\section{Ground-truth subspace and eigenvalue errors}
\label{sec:gt}
%==================================================================

Here both the real and generated samples are benchmarked against the analytic
ground-truth return covariance $\Sigma\in\R^{d\times d}$, which is known in closed
form:
\begin{equation}
  \Sigma \;=\;
  \begin{cases}
    \beta\beta\T, & \text{descriptor model, fixed } U\ (\beta=UT),\\[4pt]
    \dfrac1n\sum_{i=1}^{n}\beta_i\beta_i\T,\ \ \beta_i = U_i T,
      & \text{descriptor model, time-varying } U_i,\\[8pt]
    A_{\mathrm{load}}\T A_{\mathrm{load}} + \sigma^2 I,
      & \text{U-Net simulation (loading matrix } A_{\mathrm{load}},\ \text{noise } \sigma^2).
  \end{cases}
\end{equation}

\paragraph{Rank-$k$ truncation.}
For any symmetric PSD covariance $C$ with eigendecomposition (equivalently SVD)
$C = \sum_{j} \lambda_j\,u_j u_j\T$, $\lambda_1\ge\lambda_2\ge\cdots\ge0$, define
the top-$k$ reconstruction and the top-$k$ eigenvalue vector
\begin{equation}
  C_{(k)} = \sum_{j=1}^{k}\lambda_j\,u_j u_j\T
          = U_k\,\diag(\lambda_1,\dots,\lambda_k)\,U_k\T,
  \qquad
  \bm\lambda(C) = (\lambda_1,\dots,\lambda_k).
\end{equation}

\paragraph{Subspace error.}
With $\Sigma_{(k)}$ the ground-truth rank-$k$ reconstruction and
$C_{\mathrm r},C_{\mathrm g}$ the real / generated sample covariances of
Section~\ref{sec:moments}, define
\begin{equation}
  E^{\mathrm{sub}}_{\mathrm r} = \frac{\fro{(C_{\mathrm r})_{(k)} - \Sigma_{(k)}}}{\fro{\Sigma_{(k)}}},
  \qquad
  E^{\mathrm{sub}}_{\mathrm g} = \frac{\fro{(C_{\mathrm g})_{(k)} - \Sigma_{(k)}}}{\fro{\Sigma_{(k)}}},
  \qquad
  \frac{E^{\mathrm{sub}}_{\mathrm g}}{E^{\mathrm{sub}}_{\mathrm r}}.
  \label{eq:suberr}
\end{equation}

\paragraph{Eigenvalue error.}
With ground-truth eigenvalues $\bm\lambda(\Sigma)=(\lambda^{\mathrm{gt}}_1,\dots,\lambda^{\mathrm{gt}}_k)$,
define the mean relative top-$k$ eigenvalue error
\begin{equation}
  E^{\mathrm{eig}}_{\bullet}
    = \frac{1}{k}\sum_{j=1}^{k}
      \left|\frac{\lambda_j(C_\bullet)}{\lambda^{\mathrm{gt}}_j} - 1\right|,
    \qquad \bullet\in\{\mathrm r,\mathrm g\},
    \qquad
    \frac{E^{\mathrm{eig}}_{\mathrm g}}{E^{\mathrm{eig}}_{\mathrm r}}.
  \label{eq:eigerr}
\end{equation}
We additionally report the three raw spectra
$\bm\lambda(\Sigma),\ \bm\lambda(C_{\mathrm r}),\ \bm\lambda(C_{\mathrm g})$.
\emph{Interpretation:} $E^{\mathrm{sub}}$ measures how well the leading
return-covariance subspace is reproduced and $E^{\mathrm{eig}}$ how well its
scale (variance per factor) is reproduced. A ratio $<1$ means the
\emph{generated} sample is closer to the ground truth than the finite training
sample itself --- the goal for a well-trained generator.

%==================================================================
\section{Summary of reported quantities}
%==================================================================

\begin{table}[h]
\centering
\small
\begin{tabular}{@{}lll@{}}
\toprule
Symbol & Definition & Good value \\
\midrule
\multicolumn{3}{@{}l}{\emph{Subspace recovery ($A$ vs $T$), Section~\ref{sec:subspace}}}\\
$\theta_j,\ \overline\theta,\ \theta_{\max}$ & principal angles \eqref{eq:angles} & $\to 0$ \\
$d_{\mathrm G}$ & Grassmann distance \eqref{eq:grassmann} & $\to 0$ \\
$\Delta_P,\ \widehat\Delta_P$ & projector distance \eqref{eq:projector} & $\to 0$ \\
$\rho$ & Procrustes residual \eqref{eq:procrustes} & $\to 0$ \\
$\sigma(A),\sigma(T)$ & singular spectra & reference only \\
\midrule
\multicolumn{3}{@{}l}{\emph{Moment matching, Section~\ref{sec:moments}}}\\
$\twonorm{\mu_{\mathrm r}-\mu_{\mathrm g}}$ & mean discrepancy & $\to 0$ \\
$\epsilon_{\mathrm{cov}}$ & relative covariance error & $\to 0$ \\
$\kappa_{\mathrm{cov}},\ \kappa_{\mathrm{std}}$ & covariance / std ratio & $\to 1$ \\
\midrule
\multicolumn{3}{@{}l}{\emph{Ground-truth errors, Section~\ref{sec:gt}}}\\
$E^{\mathrm{sub}}_{\mathrm r},E^{\mathrm{sub}}_{\mathrm g}$ & subspace error \eqref{eq:suberr} & $\to 0$ \\
$E^{\mathrm{eig}}_{\mathrm r},E^{\mathrm{eig}}_{\mathrm g}$ & eigenvalue error \eqref{eq:eigerr} & $\to 0$ \\
$E^{\mathrm{sub}}_{\mathrm g}/E^{\mathrm{sub}}_{\mathrm r},\ E^{\mathrm{eig}}_{\mathrm g}/E^{\mathrm{eig}}_{\mathrm r}$ & generated-vs-training ratios & $<1$ \\
\bottomrule
\end{tabular}
\end{table}

\end{document}
